\minitopic{研读实验室：让模态判断可检验}\index{模态认识论}\index{彩票问题}\index{德性认识论}\index{鲁棒性} “在可能世界中容易出错”如果不加约束，几乎可以支持任何结论。对任一真信念，我们都能想象一个遥远世界使它为假；对任一坏方法，也能描述一个世界让它永远正确。模态认识论的任务不是比拼想象力，而是说明哪些变化与当前信念形成具有解释相关性。 一种实用做法是先固定四项：目标命题、实际形成方法、主体能力状态、背景环境；然后只改变一项，观察信念是否随事实响应。评估温度计信念时，若目标命题是“液体为20度”，形成方法是读取仪表，最近变化应包括液体温度与仪表故障，而不是主体突然获得心灵感应。若把方法改为猜测，就不再评价同一信念基础。 世界接近性还具有方向性。敏感性从$p$为假的世界寻找主体反应；安全性从主体同样相信$p$的世界检查真值。前者关心追踪，后者关心避免错误。两种条件对彩票、闭合和必要真理给出不同压力，不能用一个模糊的“反事实可靠”合并。

\minitopic{敏感性的四步检验}第一，写出$p$为假的最近方式。第二，保持方法尽量不变。第三，问主体是否仍形成对$p$的信念。第四，说明结果是否来自方法对事实的响应。例如，门禁屏幕显示“已锁”。若门未锁，传感器会显示“未锁”，相信门锁是敏感的；若屏幕只是循环播放“已锁”，则不敏感。 必要真理暴露敏感性的困难。若$p$不可能为假，就没有$p$为假的世界；把条件理解成空真，会让任何关于必要真理的信念都敏感，包括随意猜测。理论需要另行要求适当方法，或把敏感性只作为经验知识的局部条件。 闭合也制造反直觉。主体通过正常视觉敏感地知道“这是斑马”，再推出“这不是伪装骡子”。在“它是伪装骡子”的最近世界，主体可能仍相信不是伪装骡子，故结论不敏感。于是敏感性可能允许知道斑马，却不允许知道其已认识的逻辑后果。接受者可以拒绝无条件闭合，但必须说明日常演绎知识为何仍大体保留。

\minitopic{安全性的四步检验}第一，固定主体以相关方式相信$p$。第二，考察形成条件的小变化。第三，检查这些世界中$p$是否仍真。第四，排除只因重新描述方法而获得的虚假安全。停止的钟在实际世界恰显八点；几分钟前或几分钟后，主体仍以看钟方式相信相应时间，却会错，因此信念不安全。 方法具体到何种程度是争点。若描述为“恰在八点看一只停在八点的钟”，信念在所有符合描述的世界都真，但真值已被偷偷写入方法。合理类型应由主体可重复的心理或技术机制确定，并能在事前识别。安全性与可靠主义一样需要非临时的过程个体化。 绝对安全也可能排除极小故障风险。正常视觉在某个非常接近的神经突发世界会错；一台检验设备也可能罕见失灵。理论通常通过“相关近邻”或“容易”容纳可错性，但这两个词不能只按结论调节。可考虑错误机制是否已经存在、变化是否小、主体有无异常迹象，以及任务领域如何划定。

\minitopic{彩票、统计与个案联系}一张一亿分之一中奖率的彩票几乎一定会输。开奖前，主体基于概率相信“我的票会输”。这项信念很可能真，也可在许多近邻世界为真；为何不少人仍拒绝称其为知识？一种解释是，纯统计证据没有提供把这张票与失败事实连接的个案信息。每张票都得到同样理由，其中必有一张理由导向假信念。 但统计证据并非从不产生知识。医学、人口研究和质量控制大量依赖统计。差异可能不在“统计对个案”，而在证据结构：抽样是否校准，主体是否知道生成机制，结论是概率陈述还是确定个案属性，是否存在可进一步取得的决定性证据。相信“这张票中奖概率是一亿分之一”可成为知识，与相信“它必定输”不同。 法律中的“裸统计”同样敏感。若录像证明一百人中九十九人参与破坏，却无法辨认被告，仅凭比例判定该个体参与，可能在证据规范上不足。认识论要区分理性置信、知识与允许强制行动的证明标准。一个高概率信念可以理性，却不足以承担剥夺权利的制度后果。 开奖后查看官方号码，形成方法改变。此时若票中奖，号码核对会使主体不相信会输；信念敏感且安全，也获得个案联系。前后命题相同，证据基础不同，认识地位随之变化。

\minitopic{成就、功劳与合作知识}德性认识论把知识看成一种因能力而取得的真信念。射箭类比有帮助：命中若主要由风把偏离箭吹回靶心，就不因射手技术；若技术在正常条件下稳定导向命中，则成功可归功于能力。认识案例中，感知、记忆、推理与信息选择都可能是能力。 “归功”不能等于主体独立完成。研究者依赖仪器、同事与数据库仍可知道；读者从可靠书本学习也不是偶然寄生。更合适的问题是，主体是否胜任地参与一条对真相有响应的合作链：选择适当来源、理解输出、注意击败信号并在权限范围内复核。功劳可以分布，而知识归属仍有条件。 团队知识提出三种归属。第一，所有成员各自知道；第二，至少某些成员知道且信息可传递；第三，组织通过角色、记录与决策程序知道，即使无人掌握全部。若一份风险报告躺在数据库但决策流程无法调用，说“公司知道风险”可能只表达信息存储，不足以归责。能力理论必须明确所评价的主体层级。 反思能力也不是背诵机制。实验员不必理解传感器全部物理原理，但应知道校准范围、异常提示和何时求助。对高风险结论，二阶监控使一阶能力在合适环境中运行。反思要求若无限提高会排除日常知识，若完全取消又难以说明明显忽视警报为何破坏知识。

\minitopic{鲁棒性：把可能世界变成设计原则}模态条件可以转译为系统测试。敏感性对应干预：当目标事实变为假时，信息渠道是否改变输出？安全性对应扰动：当噪声、人员或输入略变时，同一输出是否仍大体为真？能力归因对应机制解释：成功是否来自设计功能和胜任使用，而非故障抵消？ 异质冗余比重复同一来源更有价值。三个使用同一数据库的模型可能共同出错；传感器、人工观察与独立历史记录响应机制不同，更能发现失败。冗余也不是越多越好，若验证通道都从原结论复制，数量只制造虚假确信。 故障注入是一种反事实实验。工程团队主动使传感器离线、改变时间戳或加入异常样本，观察系统是否报警。哲学思想实验与工程测试在此连接：前者说明应改变什么变量，后者检验真实机制如何响应。实际频率不能穷尽模态安全，却能约束哪些世界确实接近。 制度同样可以鲁棒。新闻编辑采用来源独立性检查、医学实验预注册、法院允许交叉质询，都是让结论对局部错误更不敏感的设计。它们不保证永不出错，而是让错误更可能留下可见痕迹并被纠正。

\minitopic{综合案例：洪水预警系统}城市系统结合上游传感器、气象预测和人工巡查发布预警。某次传感器因泥沙高估水位，气象模型又低估降雨，两项误差抵消，最终准确发布三级预警。居民据此相信两小时后道路会积水，事实确实如此。 总体预测正确不等于信念路径安全。若任一误差稍小，等级就错；在邻近运行中同样输出容易为假。敏感性则问，若道路不会积水，系统是否会降低预警；这取决于模型和阈值是否随事实代理变量响应。系统可能敏感却不安全，也可能在狭窄范围安全却对新型洪水不敏感。 居民知识还取决于信息链。他们无须亲自审计模型，但若官方渠道有稳定校验、明确适用范围且当前无异常警报，依赖可以是胜任的。工程师若已看到传感器冲突却压下警告，即使预警碰巧准确，其个体信念与机构责任评价会改变。 改进包括保留各模型原始输出、用独立水尺核验、对误差抵消设异常规则、发布置信区间与更新时间。措施分别改善可追踪性、个案联系、击败者可见性和使用者反思。反运气理论由此不只是给“知道”贴标签，也指导怎样建造更可知的世界。

\begin{tcolorbox}[breakable,colback=white,colframe=blue!30!black,title={本章研读任务},fonttitle=\bfseries]
选择一项真实的信息系统，画出目标事实、传感或证言渠道、主体信念与行动。分别设计一个敏感性干预、一个安全性扰动和一个错误抵消案例；再提出两种机制不同的核验方式，说明它们如何改变最近错误情形。
\end{tcolorbox}
